% Source: https://artofproblemsolving.com/community/c3h1751255_two_easy_math_problems
% Stable identity: aops-topic-1751255-q02
Triangles $AEF$ and $ABC$ share the angle at $A$, so their areas are in the ratio of the products of the sides enclosing that angle:
\[
\dfrac{[AEF]}{[ABC]}=\dfrac{\dfrac{1}{2}\,AE\cdot AF\sin A}{\dfrac{1}{2}\,AB\cdot AC\sin A}=\dfrac{AE}{AB}\cdot\dfrac{AF}{AC}=\dfrac{1}{3}\cdot\dfrac{1}{2}=\dfrac{1}{6} .
\]
(Without trigonometry: $[AEC]=\dfrac{1}{3}[ABC]$ because the two triangles have the same height from $C$ and bases $AE=\dfrac{1}{3}AB$; and $[AEF]=\dfrac{1}{2}[AEC]$ because $F$ is the midpoint of $AC$.)

\begin{asy}
usepackage("amsmath");
size(6cm);
defaultpen(fontsize(10pt));
pair A=(1.6,5),B=(0,0),C=(7,0);
pair Ep=A+(B-A)/3,F=(A+C)/2;
fill(A--Ep--F--cycle,palegray);
draw(A--B--C--cycle,linewidth(0.9));
draw(Ep--F,linewidth(0.8));
dot(Ep);dot(F);
label("$A$",A,N);label("$B$",B,SW);label("$C$",C,SE);
label("$E$",Ep,W);label("$F$",F,NE);
label("$1$",(A+Ep)/2,W);label("$2$",(Ep+B)/2,W);
label("$1$",(A+F)/2,NE);label("$1$",(F+C)/2,NE);
label("\mbox{}",(min(currentpicture,user=true)+max(currentpicture,user=true))/2,black); // keeps the text after the figure black when LaTeX typesets the labels inline
\end{asy}

Hence $[AEF]=\dfrac{72}{6}=\boxed{12}$.
